\documentclass[11pt]{article}
\usepackage[margin=1in]{geometry}
\usepackage{booktabs}
\usepackage{graphicx}
\usepackage{amsmath}
\usepackage{amssymb}
\usepackage{siunitx}
\usepackage{hyperref}
\usepackage{natbib}
\usepackage{xcolor}
\usepackage{authblk}

\title{Chronic Activation of a Negative Engram Induces Behavioral and Glial Changes in a Mouse Model of Aging}
\author{Anonymous}
\date{}

\begin{document}
\maketitle

\begin{abstract}
Persistent recall of aversive experiences is a core feature of anxiety disorders, post-traumatic stress, and age-associated cognitive decline, yet the causal impact of chronic reactivation of a negative memory engram on the aging brain is poorly understood. We used the TRAP2 (Fos--iCreER$^{\mathrm{T2}}$) mouse line combined with Cre-dependent excitatory designer receptors exclusively activated by designer drugs (DREADDs; AAV9-DIO-Flex-hM3Dq-mCherry) to chemogenetically tag and chronically reactivate a contextual fear-conditioned engram in ventral hippocampal CA1 (vCA1) of young (3-month) and old (11-month) male mice across three months of continuous deschloroclozapine (DCZ, 3\,\si{\micro\gram\per\kilogram\per\day} in home-cage drinking water) administration. After this stimulation period, 6- and 14-month-old hM3Dq mice spent significantly less time in the center of the open field (Two-way ANOVA, Group effect $F(1,33)=29.56$, $p<0.0001$) and the open arms of the zero maze (Group $F(1,38)=40.22$, $p<0.0001$) relative to mCherry controls. Working memory, assessed by spontaneous alternation in the Y-maze, was selectively impaired in the 14-month hM3Dq group (Tukey's post-hoc $p=0.0045$). Chronic engram reactivation impaired extinction of remote fear memory in 6-month hM3Dq mice (Group $F(1,17)=5.494$, $p=0.0315$) and drove generalization of freezing to a novel context B at both ages (Group $F(1,41)=26.53$, $p<0.0001$). Histological quantification revealed that chronic engram activation did not reduce NeuN$^{+}$ cell counts in subiculum, ventral dentate gyrus, or vCA1, but increased the number and altered the morphology of Iba1$^{+}$ microglia (14-month only) and GFAP$^{+}$ astrocytes (both ages; Group $F(1,8)=26.95$, $p=0.0008$). Collectively, three months of vCA1 negative-engram activation is sufficient to recapitulate anxiety-like, fear-generalization, and glial phenotypes reminiscent of early-stage age-related pathology, without overt neurotoxicity.
\end{abstract}

\section{Introduction}
The hippocampus supports the encoding, consolidation, and contextual recall of episodic and emotionally salient memories. In rodents, individual memories can be read out at the level of sparse ensembles --- engrams --- that are reliably tagged by activity-dependent immediate-early gene expression. Chemogenetic manipulation of engram-defined neurons thus provides a causal handle on how specific memory traces shape cognition and affect over time. A growing literature implicates sustained or inappropriate reactivation of negatively valenced engrams in the etiology of anxiety, depression, and post-traumatic stress; yet almost all published engram manipulations span hours to days, leaving open the question of how chronic, months-long drive of a negative engram remodels behavior and the cellular milieu of the hippocampus, particularly against the backdrop of normal aging.

Aging is associated with progressive hippocampal vulnerability and cumulative risk for neurodegenerative disease. \emph{Apoe} genotype accounts for the majority of Alzheimer's disease (AD) risk and pathology \citep{johnson2015,schuff2009}, with the E4 isotype accounting for as much as 50\% of AD in the United States \citep{raber2003}. Whether prolonged reactivation of a maladaptive memory trace --- a cognitive substrate of rumination --- can itself accelerate hippocampal dysfunction and glial reactivity during middle age is unknown.

Here we use TRAP2 mice \citep{denardo2019,monasterio2024} in combination with Cre-dependent hM3Dq DREADDs targeted to vCA1 pyramidal neurons to tag a contextual fear-conditioning (CFC) engram and chronically reactivate it for three months via home-cage delivery of water-soluble deschloroclozapine (DCZ; \citealt{nagai2020,suthard2023}). Young (3-month) and old (11-month) cohorts were then subjected to a battery of assays probing anxiety (open field, zero maze), working memory (Y-maze), remote recall, extinction, and generalization of fear, followed by quantitative histology of neurons, microglia, and astrocytes in the ventral hippocampus. We report that three months of negative engram activation is sufficient to induce anxiety-like behavior and context generalization at both ages, age-selective impairments of working memory and extinction, and increases in number and morphological complexity of hippocampal glia --- all in the absence of detectable neuronal loss.

\section{Methods}

\subsection{Animals and ethics}
Experiments used male TRAP2 (Fos--iCreER$^{\mathrm{T2}}$) mice in two age cohorts --- young (3 months at tagging, 6 months at behavioral testing) and old (11 months at tagging, 14 months at testing) --- randomly assigned to receive AAV9-DIO-Flex-hM3Dq-mCherry (hM3Dq) or AAV9-DIO-Flex-mCherry (control). Following surgery, mice were group-housed with littermates and allowed to recover for a minimum of 4 weeks before experimentation to permit viral expression. All subjects were treated in accord with protocol 201800579 approved by the Institutional Animal Care and Use Committee (IACUC) at Boston University.

\subsection{Stereotaxic surgery}
Mice were initially anesthetized with 3.0\% isoflurane during induction and maintained at 1--2\% isoflurane through a stereotaxic nose-cone delivery (oxygen 1\,L/min). Hair above the surgical site was removed using Veet cream and the skin was cleaned with alternating betadine and 70\% ethanol. 2.0\% lidocaine hydrochloride was injected subcutaneously as local analgesia prior to midsagittal incision. A 0.1\,mg/kg (5\,mg/kg stock) subcutaneous dose of meloxicam was administered at the beginning of surgery. All animals received bilateral craniotomies with a 0.5--0.6\,mm drill-bit for vCA1 injections. A 10\,\si{\micro\liter} airtight Hamilton syringe with a 33-gauge beveled needle was slowly lowered to vCA1: $-3.16$\,mm anteroposterior (AP), $\pm 3.10$\,mm mediolateral (ML), and both $-4.25$ and $-4.50$\,mm dorsoventral (DV) to cover the vertical axis of the ventral hippocampus. A total volume of 300\,nL of AAV9-DIO-Flex-hM3Dq-mCherry or AAV9-DIO-Flex-mCherry was bilaterally injected into vCA1 (150\,nL per DV site). Following surgery, mice received a 0.1\,mg/kg intraperitoneal (IP) dose of buprenorphine, scaled by body weight. Histological assessment verified bilateral viral targeting; data from off-target injections were excluded from analyses.

\subsection{Engram tagging and chronic DREADD activation}
Two weeks after surgery, all mice underwent CFC consisting of four 1.5\,mA foot shocks in Context A. 30 minutes after CFC, mice received a 40\,mg/kg injection of 4-OHT (4-hydroxytamoxifen) to induce recombination and permanently tag cells active during the fearful experience. All animals then had their home-cage drinking water replaced every 5 days with water-soluble deschloroclozapine dihydrochloride (DCZ; HelloBio \#HB9126) in distilled water (diH$_2$O). A stock was prepared by dissolving 1--2\,mg DCZ in 1\,mL DMSO and vortexing; 17\,\si{\micro\liter} of stock was combined with 1\,L of diH$_2$O, yielding a final in-cage concentration calibrated to deliver $\sim$3\,\si{\micro\gram\per\kilogram\per\day} assuming $\sim$6\,mL of water per mouse per day, following \citet{suthard2023}. DCZ peaks 5--10 minutes after IP administration and declines to baseline by $\sim$2\,h \citep{nagai2020}. No signs of distress or seizures were observed over the entire three-month administration.

\subsection{Behavioral assays}
After three months of stimulation (young cohort now 6 months; old cohort now 14 months), mice were handled for 3--5 days (2\,min/day) before behavioral testing. The battery included open field, Y-maze, zero maze, remote recall, extinction, and generalization to a novel Context B (Fig.~2A). The open field arena measured 61$\times$61\,cm with black plastic walls and a red-taped central zone; animals explored for 10\,min. The zero maze consisted of a 44\,cm-wide circular ring with outer diameter 62\,cm, 8\,cm in height, elevated 66\,cm off the floor. The Y-maze consisted of clear plexiglass arms (36.5\,cm length, 7.5\,cm width, 12.5\,cm height) intersected at 120\textdegree. Fear conditioning was performed in a 18.5$\times$18.5$\times$21.5\,cm chamber with metal walls, plexiglass front/back, and a 16-bar stainless-steel grid floor. Chambers were cleaned with 70\% ethanol between animals. For generalization (Context B), 5 days after extinction, animals were taken to a different room and placed in a new, larger chamber. Remote recall consisted of re-exposure to Context A for 5\,min (300\,s) without shocks. Extinction consisted of 30\,min (1800\,s) in Context A without shocks.

\subsection{Histology and imaging}
Mice were overdosed with 3\% isoflurane and perfused transcardially with cold (4\,\textdegree C) PBS followed by 4\% paraformaldehyde (PFA) in PBS. Brains were post-fixed in PFA at 4\,\textdegree C for 24--48\,h, cryoprotected in 30\% sucrose/PBS, and stored in 0.01\% sodium azide/PBS until sectioning. Brains were cut at 50\,\si{\micro\meter} on a vibratome into coronal and sagittal sections. Sections were washed three times for 10\,min in PBS, blocked for 2\,h at room temperature in PBS with 0.2\% Triton X-100 (PBST) and 5\% normal goat serum, then incubated with primary and secondary antibodies against NeuN, Iba1, and GFAP. NeuN z-stacks were acquired at 4\,\si{\micro\meter} intervals; Iba1 and GFAP z-stacks at 1\,\si{\micro\meter} intervals to support detailed morphological analysis. Morphology of Iba1$^{+}$ microglia and GFAP$^{+}$ astrocytes was quantified using the MATLAB-based tool 3DMorph.

\subsection{RNA-seq (FACS-purified vCA1 populations)}
FACS-sorted cells were passed through a MACS SmartStrainer (70\,\si{\micro\meter}) before library preparation. RNA-seq libraries were prepared using the SMART-Seq v4 Ultra Low Input RNA Kit (TaKaRa Bio, USA). For Gene Ontology analysis, protein-coding genes were ranked by fold change and used as input for the GseaPreranked tool.

\subsection{Statistics}
Data were analyzed in GraphPad Prism v9.4.1 (La Jolla, CA). Two-way ANOVA with Tukey's multiple comparisons test was used for between-subject designs; two-way repeated-measures (RM) ANOVA with Sidak's post-hoc test was used for within-subject designs. Outliers were identified with ROUT ($Q=10\%$, increased FDR to raise power). Sample sizes ranged from $n=8$--17 per group after outlier removal for behavioral endpoints and $n=3$ mice/group $\times$ 18 single-tile ROIs for cell counts (NeuN, Iba1, GFAP); morphology used $n=3$ mice/group $\times$ 3 single-tile ROIs, each containing 20--30 glial cells. Significance was set at $^{*}p\leq 0.05$, $^{**}p\leq 0.01$, $^{***}p\leq 0.001$, $^{****}p\leq 0.0001$.

Table~\ref{tab:design} summarizes the experimental design.

\begin{table}[h]
\centering
\caption{Experimental design. Two cohorts (young, old) were each split between hM3Dq and mCherry viruses, yielding four groups. Ages refer to tagging (3MO/11MO) and behavior (6MO/14MO) time points.}
\label{tab:design}
\begin{tabular}{lllll}
\toprule
Group & Virus (vCA1) & Age at tagging & Age at testing & DCZ duration \\
\midrule
6MO mCherry & AAV9-DIO-Flex-mCherry   & 3 months  & 6 months  & 3 months \\
6MO hM3Dq   & AAV9-DIO-Flex-hM3Dq-mCherry & 3 months & 6 months  & 3 months \\
14MO mCherry & AAV9-DIO-Flex-mCherry  & 11 months & 14 months & 3 months \\
14MO hM3Dq   & AAV9-DIO-Flex-hM3Dq-mCherry & 11 months & 14 months & 3 months \\
\bottomrule
\end{tabular}
\end{table}

\section{Results}

\subsection{Chronic hM3Dq activation does not disrupt normal weight gain or induce neuron loss}
Body weights were recorded monthly across the three-month stimulation (Fig.~2C). Both young and old mice gained weight independently of hM3Dq activation. For the 6-month cohort (Two-way RM ANOVA), Interaction $F(2,34)=0.2633$, $p=0.7701$; Month $F(2,34)=14.84$, $p<0.0001$; Group $F(1,17)=1.361$, $p=0.2594$; Subject $F(17,34)=15.44$, $p<0.0001$ (Sidak's Month: 1 vs. 2, $p=0.0279$; 1 vs. 3, $p<0.0001$; 2 vs. 3, $p=0.0323$). For the 14-month cohort, Interaction $F(2,52)=2.607$, $p=0.0833$; Month $F(2,52)=7.854$, $p=0.0010$; Group $F(1,26)=0.1275$, $p=0.7239$; Subject $F(26,52)=20.74$, $p<0.0001$ (Sidak's Month: 1 vs. 2, $p=0.6942$; 1 vs. 3, $p=0.0011$; 2 vs. 3, $p=0.0198$). We note the absence of a Month 0 baseline weight measurement; it is therefore possible that body weight dipped and rebounded during Months 0--1.

Importantly, three months of chronic activation did not reduce the number of NeuN$^{+}$ cells in subiculum, ventral dentate gyrus (vDG), or vCA1. Two-way ANOVA for subiculum: Interaction $F(1,8)=0.2516$, $p=0.6295$; Age $F(1,8)=7.341$, $p=0.0267$; Group $F(1,8)=0.00721$, $p=0.9344$ (Tukey's Group: 6MO mCherry vs. hM3Dq $p=0.9904$; 14MO mCherry vs. hM3Dq $p=0.9744$). For vDG: Interaction $F(1,8)=0.6806$, $p=0.4333$; Age $F(1,8)=4.423$, $p=0.0686$; Group $F(1,8)=0.0524$, $p=0.8246$. For vCA1: Interaction $F(1,8)=2.824$, $p=0.1314$; Age $F(1,8)=0.5443$, $p=0.4817$; Group $F(1,8)=2.947$, $p=0.1244$. Thus DCZ-driven chronic engram activation does not produce gross neuronal loss in the ventral hippocampus.

\subsection{Chronic negative-engram activation increases anxiety-like behavior}
In the open field (Fig.~3A), 6- and 14-month-old hM3Dq mice spent less time in the center compared to age-matched mCherry controls (Two-way ANOVA; Interaction $F(1,33)=0.1438$, $p=0.7070$; Age $F(1,33)=0.5858$, $p=0.4495$; Group $F(1,33)=29.56$, $p<0.0001$; Tukey's Group: 6MO mCherry vs. hM3Dq $p=0.0049$; 14MO mCherry vs. hM3Dq $p=0.0016$). No significant differences were detected in total distance traveled (Interaction $F(1,37)=0.0016$, $p=0.9675$; Age $F(1,37)=0.6862$, $p=0.4128$; Group $F(1,37)=0.3522$, $p=0.5565$), nor in the number of center entries (Interaction $F(1,41)=1.236$, $p=0.2728$; Age $F(1,41)=0.4424$, $p=0.5097$; Group $F(1,41)=0.3370$, $p=0.5648$), indicating that mCherry and hM3Dq mice enter the center equally often but spend less time per visit. 6-month-old hM3Dq mice additionally showed a higher average speed (Interaction $F(1,38)=4.003$, $p=0.526$; Age $F(1,38)=0.2346$, $p=0.6306$; Group $F(1,38)=7.306$, $p=0.0102$; Tukey's Group: 6MO mCherry vs. hM3Dq $p=0.0150$; 14MO mCherry vs. hM3Dq $p=0.9532$).

In the zero maze (Fig.~3B), both hM3Dq groups spent less time in the open arms relative to age-matched controls (Interaction $F(1,38)=0.5633$, $p=0.4575$; Age $F(1,38)=0.1709$, $p=0.6816$; Group $F(1,38)=40.22$, $p<0.0001$; Tukey's Group: 6MO mCherry vs. hM3Dq $p=0.0002$; 14MO mCherry vs. hM3Dq $p=0.0007$). No differences were observed in total distance (Interaction $F(1,36)=1.915$, $p=0.1749$; Age $F(1,36)=0.0093$, $p=0.9237$; Group $F(1,36)=0.6022$, $p=0.4428$), number of open-area entries (Interaction $F(1,38)=1.082$, $p=0.3049$; Age $F(1,38)=0.3839$, $p=0.5392$; Group $F(1,38)=3.286$, $p=0.0778$), or average speed (Interaction $F(1,37)=0.5118$, $p=0.4788$; Age $F(1,37)=1.659$, $p=0.2057$; Group $F(1,37)=0.7180$, $p=0.4023$). These results indicate a robust, age-independent anxiogenic effect of chronic negative-engram activation. Key effects are summarized in Table~\ref{tab:anxiety}.

\begin{table}[h]
\centering
\caption{Chronic vCA1 engram activation produces anxiety-like behavior in open field and zero maze. Two-way ANOVA main Group effects and age-matched Tukey's post-hoc comparisons.}
\label{tab:anxiety}
\begin{tabular}{llcc}
\toprule
Assay (endpoint) & ANOVA Group effect & 6MO mC vs. hM3Dq & 14MO mC vs. hM3Dq \\
\midrule
Open field --- time in center   & $F(1,33)=29.56$, $p<0.0001$ & $p=0.0049$ & $p=0.0016$ \\
Open field --- distance         & $F(1,37)=0.3522$, $p=0.5565$ & ns & ns \\
Open field --- center entries   & $F(1,41)=0.3370$, $p=0.5648$ & ns & ns \\
Open field --- avg speed        & $F(1,38)=7.306$, $p=0.0102$  & $p=0.0150$ & $p=0.9532$ \\
Zero maze --- time in open arms & $F(1,38)=40.22$, $p<0.0001$ & $p=0.0002$ & $p=0.0007$ \\
Zero maze --- distance          & $F(1,36)=0.6022$, $p=0.4428$ & ns & ns \\
Zero maze --- open-area entries & $F(1,38)=3.286$, $p=0.0778$  & ns & ns \\
Zero maze --- avg speed         & $F(1,37)=0.7180$, $p=0.4023$ & ns & ns \\
\bottomrule
\end{tabular}
\end{table}

\subsection{Working memory is impaired only in aged mice}
In the Y-maze (Fig.~3C), 14-month hM3Dq mice showed a lower percentage of spontaneous alternations than 14-month mCherry mice (Two-way ANOVA; Interaction $F(1,40)=4.088$, $p=0.0499$; Age $F(1,40)=0.3080$, $p=0.5820$; Group $F(1,40)=7.517$, $p=0.0091$; Tukey's Group: 14MO mCherry vs. hM3Dq $p=0.0045$). No difference was observed in the 6-month cohort (Tukey's Group: 6MO mCherry vs. hM3Dq $p=0.9634$). Neither total distance traveled (Interaction $F(1,44)=0.2668$, $p=0.6081$; Age $F(1,44)=0.5106$, $p=0.4787$; Group $F(1,44)=1.181$, $p=0.2831$) nor total number of arm entries (Interaction $F(1,40)=1.332$, $p=0.2553$; Age $F(1,40)=0.2236$, $p=0.6389$; Group $F(1,40)=0.8247$, $p=0.3692$) differed across groups, indicating a specific deficit in spatial working memory rather than overall locomotion.

\subsection{Negative-engram activation impairs extinction and produces fear generalization}
During CFC, freezing across the 300\,s session did not differ between hM3Dq and mCherry within either age group, as expected (Two-way ANOVA across total freezing; Interaction $F(1,42)=0.1342$, $p=0.7150$; Age $F(1,42)=12.20$, $p=0.0011$; Group $F(1,42)=2.913$, $p=0.0953$). RM analyses of freezing across time further confirmed robust acquisition (6MO: Time $F(5,130)=115.6$, $p<0.0001$; 14MO: Time $F(5,90)=234.4$, $p<0.0001$) with no Group effect (6MO: Group $F(1,26)=2.737$, $p=0.1101$; 14MO: Group $F(1,18)=1.862$, $p=0.1892$).

Three months later, mice were returned to Context A for remote recall (300\,s; Fig.~4B). In the 6-month cohort, RM ANOVA revealed a Group effect ($F(1,17)=6.588$, $p=0.0200$) and Interaction ($F(5,85)=15.72$, $p=0.0115$), with 6MO hM3Dq mice freezing more than 6MO mCherry; in the 14-month cohort, no Group effect was detected ($F(1,26)=0.3967$, $p=0.5343$; Interaction $F(5,130)=0.9027$, $p=0.4815$; Time $F(5,130)=14.18$, $p<0.0001$). Aggregated total freezing during remote recall confirmed a Group effect (Two-way ANOVA; Interaction $F(1,41)=2.531$, $p=0.1193$; Age $F(1,41)=10.08$, $p=0.0028$; Group $F(1,41)=6.406$, $p=0.0153$; Tukey's Age: 6 vs. 14MO mCherry $p=0.0101$, 6 vs. 14MO hM3Dq $p=0.6672$; Tukey's Group: 6MO mCherry vs. hM3Dq $p=0.0449$; 14MO mCherry vs. hM3Dq $p=0.8890$).

Extinction over 30\,min (1800\,s; Fig.~4C) in Context A revealed an impaired ability of 6MO hM3Dq mice to extinguish fear (RM ANOVA 6MO: Interaction $F(5,85)=0.4133$, $p=0.8383$; Time $F(5,85)=1.979$, $p=0.0899$; Group $F(1,17)=5.494$, $p=0.0315$; Subject $F(17,85)=7.352$, $p<0.0001$). In 14MO mice, extinction was comparable between groups (Interaction $F(5,130)=1.362$, $p=0.2570$; Time $F(5,130)=5.387$, $p=0.0002$; Group $F(1,26)=0.3836$, $p=0.5411$; Subject $F(26,130)=11.16$, $p<0.0001$). Follow-up analysis of the 300\,s generalization test in Context B (Fig.~4D) showed a trend toward increased freezing in 6-month hM3Dq mice (Two-way ANOVA; Interaction $F(1,41)=2.618$, $p=0.1133$; Age $F(1,41)=2.335$, $p=0.1342$; Group $F(1,41)=5.763$, $p=0.0210$; Tukey's Group: 6MO mCherry vs. hM3Dq $p=0.0590$; 14MO mCherry vs. hM3Dq $p=0.9269$).

Strikingly, when total freezing in Context B was considered across both age groups, hM3Dq mice of both ages froze more in a neutral context than age-matched mCherry controls (Two-way ANOVA; Interaction $F(1,41)=0.09801$, $p=0.7558$; Age $F(1,41)=2.624$, $p=0.1129$; Group $F(1,41)=26.53$, $p<0.0001$; Tukey's Group: 6MO mCherry vs. hM3Dq $p=0.0053$; 14MO mCherry vs. hM3Dq $p=0.0026$), indicating pronounced fear generalization. RM ANOVA in 14MO mice confirmed a Group effect ($F(1,25)=6.343$, $p=0.0186$) and Time effect ($F(4,100)=5.717$, $p=0.0003$), while 6MO mice showed a trend ($F(1,18)=4.151$, $p=0.0566$). Table~\ref{tab:fear} summarizes the fear-memory phenotypes.

\begin{table}[h]
\centering
\caption{Fear-memory phenotypes following three months of chronic vCA1 engram activation. Reported statistics are from the two-way ANOVA of total freezing per session, except where a RM ANOVA Group effect is explicitly called out.}
\label{tab:fear}
\begin{tabular}{llc}
\toprule
Phase (duration) & Statistic & Tukey's age-matched posthoc \\
\midrule
CFC (300\,s)          & Group $F(1,42)=2.913$, $p=0.0953$ & ns (acquisition matched) \\
Remote recall total  & Group $F(1,41)=6.406$, $p=0.0153$ & 6MO $p=0.0449$; 14MO $p=0.8890$ \\
Remote recall 6MO RM & Group $F(1,17)=6.588$, $p=0.0200$ & --- \\
Extinction 6MO RM    & Group $F(1,17)=5.494$, $p=0.0315$ & --- \\
Extinction 14MO RM   & Group $F(1,26)=0.3836$, $p=0.5411$ & --- \\
Generalization (Context B) total & Group $F(1,41)=26.53$, $p<0.0001$ & 6MO $p=0.0053$; 14MO $p=0.0026$ \\
Generalization 14MO RM & Group $F(1,25)=6.343$, $p=0.0186$ & --- \\
\bottomrule
\end{tabular}
\end{table}

\subsection{Chronic engram activation remodels hippocampal glia without neuronal loss}
For microglia (Fig.~5A--J), 14-month hM3Dq mice showed an increased number of Iba1$^{+}$ cells in the hippocampus (Two-way ANOVA; Interaction $F(1,8)=7.615$, $p=0.0247$; Age $F(1,8)=3.259$, $p=0.1087$; Group $F(1,8)=13.42$, $p=0.0064$; Tukey's Group: 6MO mCherry vs. hM3Dq $p=0.9906$; 14MO mCherry vs. hM3Dq $p=0.0113$; Tukey's Age: 6 vs. 14MO mCherry $p=0.9876$, 6 vs. 14MO hM3Dq $p=0.0704$). In the same 14MO hM3Dq animals, Iba1$^{+}$ microglia displayed increased ramification index (Interaction $F(1,8)=1.908$, $p=0.2046$; Age $F(1,8)=1.950$, $p=0.2001$; Group $F(1,8)=15.11$, $p=0.0046$; Tukey's 14MO mCherry vs. hM3Dq $p=0.03242$), increased cell-territory volume (Group $F(1,8)=16.45$, $p=0.0037$; 14MO mCherry vs. hM3Dq $p=0.0129$), increased minimum branch length (Group $F(1,8)=7.470$, $p=0.0257$; 14MO mCherry vs. hM3Dq $p=0.0174$), and increased maximum branch length (Group $F(1,8)=8.964$, $p=0.0172$; 14MO mCherry vs. hM3Dq $p=0.0309$). No differences were detected in average centroid distance (Group $F(1,8)=5.459$, $p=0.0477$; 14MO mCherry vs. hM3Dq $p=0.1838$), number of endpoints (Group $F(1,8)=0.0237$, $p=0.8815$), or average branch length (Group $F(1,8)=0.1211$, $p=0.7369$).

For astrocytes (Fig.~5K--T), both 6- and 14-month hM3Dq groups showed increased GFAP$^{+}$ cell counts (Two-way ANOVA; Interaction $F(1,8)=0.05526$, $p=0.8201$; Age $F(1,8)=6.824$, $p=0.0310$; Group $F(1,8)=26.95$, $p=0.0008$; Tukey's Group: 6MO mCherry vs. hM3Dq $p=0.0328$; 14MO mCherry vs. hM3Dq $p=0.0208$). Astrocytes at both ages had increased average branch length (Interaction $F(1,8)=0.1537$, $p=0.7052$; Age $F(1,8)=0.9883$, $p=0.3493$; Group $F(1,8)=32.75$, $p=0.0004$; Tukey's Group: 6MO mCherry vs. hM3Dq $p=0.0228$; 14MO mCherry vs. hM3Dq $p=0.0109$) and maximum branch length (Group $F(1,8)=41.82$, $p=0.0002$; 6MO mCherry vs. hM3Dq $p=0.0041$; 14MO mCherry vs. hM3Dq $p=0.0155$). Minimum branch length was selectively increased in 6-month hM3Dq astrocytes (Interaction $F(1,8)=1.568$, $p=0.2459$; Age $F(1,8)=0.7927$, $p=0.3993$; Group $F(1,8)=16.26$, $p=0.0038$; Tukey's Group: 6MO mCherry vs. hM3Dq $p=0.0238$; 14MO mCherry vs. hM3Dq $p=0.2759$). No differences reached significance in ramification index (Group $F(1,8)=6.677$, $p=0.0324$; 6MO mCherry vs. hM3Dq $p=0.3233$; 14MO mCherry vs. hM3Dq $p=0.3339$), cell-territory volume (Group $F(1,8)=1.996$, $p=0.1954$), average centroid distance (Group $F(1,8)=3.856$, $p=0.0852$), or number of endpoints (Group $F(1,8)=0.0071$, $p=0.9346$). Table~\ref{tab:glia} summarizes the glial phenotype.

\begin{table}[h]
\centering
\caption{Glial remodeling after three months of chronic engram activation. Two-way ANOVA Group main effects (all with $(1,8)$ numerator/denominator df) with age-matched Tukey's post-hoc comparisons.}
\label{tab:glia}
\begin{tabular}{llcc}
\toprule
Cell type / endpoint & Group $F$, $p$ & 6MO mC vs. hM3Dq & 14MO mC vs. hM3Dq \\
\midrule
Iba1$^{+}$ cell count           & $F=13.42$, $p=0.0064$ & $p=0.9906$ & $p=0.0113$ \\
Iba1$^{+}$ ramification index   & $F=15.11$, $p=0.0046$ & $p=0.3513$ & $p=0.03242$ \\
Iba1$^{+}$ cell territory volume & $F=16.45$, $p=0.0037$ & $p=0.4564$ & $p=0.0129$ \\
Iba1$^{+}$ min branch length    & $F=7.470$, $p=0.0257$ & $p=0.9996$ & $p=0.0174$ \\
Iba1$^{+}$ max branch length    & $F=8.964$, $p=0.0172$ & $p=0.8994$ & $p=0.0309$ \\
Iba1$^{+}$ avg centroid dist.   & $F=5.459$, $p=0.0477$ & $p=0.7370$ & $p=0.1838$ \\
Iba1$^{+}$ \# endpoints         & $F=0.0237$, $p=0.8815$ & ns & ns \\
Iba1$^{+}$ avg branch length    & $F=0.1211$, $p=0.7369$ & ns & ns \\
GFAP$^{+}$ cell count           & $F=26.95$, $p=0.0008$ & $p=0.0328$ & $p=0.0208$ \\
GFAP$^{+}$ avg branch length    & $F=32.75$, $p=0.0004$ & $p=0.0228$ & $p=0.0109$ \\
GFAP$^{+}$ max branch length    & $F=41.82$, $p=0.0002$ & $p=0.0041$ & $p=0.0155$ \\
GFAP$^{+}$ min branch length    & $F=16.26$, $p=0.0038$ & $p=0.0238$ & $p=0.2759$ \\
GFAP$^{+}$ ramification index   & $F=6.677$, $p=0.0324$ & $p=0.3233$ & $p=0.3339$ \\
GFAP$^{+}$ cell territory volume & $F=1.996$, $p=0.1954$ & ns & ns \\
GFAP$^{+}$ avg centroid dist.   & $F=3.856$, $p=0.0852$ & ns & ns \\
GFAP$^{+}$ \# endpoints         & $F=0.0071$, $p=0.9346$ & ns & ns \\
\bottomrule
\end{tabular}
\end{table}

\section{Discussion}
Three months of chemogenetic reactivation of a vCA1 negative-valence engram, beginning in young adulthood or middle age, is sufficient to induce a constellation of behavioral and cellular changes that parallel early features of age-associated cognitive decline. At the behavioral level, sustained negative-engram drive produced pronounced anxiety-like behavior at both ages (reduced center time in the open field and reduced open-arm time in the zero maze), impaired extinction of remote fear memory in the 6-month cohort, and broad generalization of freezing to a novel context across both ages. In aged 14-month mice, we further observed a specific impairment of spatial working memory in the Y-maze, consistent with a layered age-by-engram interaction. At the cellular level, chronic activation did not reduce NeuN$^{+}$ neuron counts in subiculum, vDG, or vCA1, indicating that the behavioral effects are unlikely to arise from overt neurodegeneration at this stimulation regime. Instead, we observed increased number and morphological remodeling of hippocampal microglia (selectively in 14-month hM3Dq mice) and astrocytes (at both ages). Collectively these results argue that chronic drive of a negative memory engram is a cognitive/affective insult capable of reshaping hippocampal glia before --- and independently of --- frank neuron loss.

Our chemogenetic strategy depends on tagging a fear-encoding ensemble using the TRAP2 Fos--iCreER$^{\mathrm{T2}}$ allele together with 4-OHT. TRAP2 labels a subset of endogenously active CA1 pyramidal cells (on the order of 5--18\%), whereas Fos--tTA/TetO-based (``DOX'') systems label 20--40\% \citep{denardo2019,monasterio2024}, and the two systems differ in their efficacy and temporal resolution based on genetic cell type, brain region, timing of 4-OHT or doxycycline delivery, metabolic variability, and sensitivity thresholds. An advantage of the present TRAP2$\times$DCZ design is the ability to maintain DREADD drive at $\sim$3\,\si{\micro\gram\per\kilogram\per\day} via drinking water for three continuous months without overt signs of distress or seizures, complementing a shorter-timescale engram literature.

Our glial results place sustained negative-engram activation within a broader neuroimmune framework for cognitive decline. While \emph{Apoe} genotype accounts for the majority of AD risk and pathology \citep{johnson2015,schuff2009} --- with the E4 isotype responsible for as much as 50\% of AD in the U.S.\ \citep{raber2003} --- our data indicate that an environmental, purely cognitive perturbation --- repeated reactivation of a negative memory --- is itself sufficient to trigger astrogliosis and, in aged brain, microglial remodeling. The accompanying behavioral profile (anxiety, extinction failure, generalization) is reminiscent of features seen in mood-disorder and early AD-risk populations, suggesting that chronic cognitive engagement of fear circuits may be a convergent driver of neuroinflammation and affective disruption.

Several caveats merit emphasis. First, chronic assessment of DREADD impact on excitatory/inhibitory (E/I) balance must eventually be quantified with both excitatory and inhibitory markers to understand how chemogenetic engram modulation reshapes circuit-level E/I dynamics. Other work has reported DREADD toxicity at high titer levels of AAV2/7-CaMKII-hM4Di-mCherry in the hippocampus at five weeks \citep{goossens2021}; our AAV9-DIO-Flex-hM3Dq-mCherry protocol, at the titer and three-month duration used here, did not reduce NeuN$^{+}$ counts in any region examined but warrants continued titer- and serotype-controlled comparisons. Second, remote recall was tested at three months post-conditioning (as opposed to a shorter 28-day window), which limits direct comparison to more standard remote-memory studies but captures a longer-timescale dysregulation of fear. Third, because a Month 0 baseline body weight was not collected, short-term weight fluctuations in the first month of DCZ administration cannot be excluded, even though longer-term weight gain was preserved. Fourth, pharmacokinetics of DCZ peak 5--10\,min after IP injection and return to baseline by $\sim$2\,h \citep{nagai2020}; continuous home-cage delivery is therefore likely to produce a pulsatile rather than tonic chemogenetic drive, the precise kinetics of which remain to be mapped \emph{in vivo}.

In summary, we establish that chronic, three-month activation of a vCA1 negative engram is sufficient to recapitulate anxiety, fear-generalization, extinction, and glial signatures reminiscent of aging-related cognitive dysfunction, without overt neurodegeneration. This paradigm provides a tractable causal model to dissect how sustained cognitive reactivation of a negative trace drives hippocampal and behavioral decline, and offers a framework for mechanistic and therapeutic studies targeting maladaptive engram persistence during aging.

\bibliographystyle{plainnat}
\bibliography{references}

\end{document}
